\documentclass[tikz,border=6pt]{standalone}
\usepackage{ctex}
\usepackage{amsmath}
\usetikzlibrary{arrows.meta,calc}

\newcommand{\psid}{\psi_d}
\newcommand{\psiref}{\psi_{\mathrm{ref}}}
\newcommand{\ud}{u_d}

\begin{document}

\begin{tikzpicture}[
  >=Stealth,
  font=\footnotesize,
  block/.style={draw, rounded corners=2pt, thick, minimum width=2.2cm, minimum height=0.9cm, align=center, fill=blue!6, draw=blue!55!black},
  core/.style={draw, rounded corners=3pt, very thick, minimum width=2.4cm, minimum height=1.0cm, align=center, fill=orange!12, draw=orange!75!black},
  plant/.style={draw, rounded corners=2pt, thick, minimum width=2.0cm, minimum height=0.9cm, align=center, fill=green!8, draw=green!45!black},
  note/.style={draw, rounded corners=2pt, align=left, fill=gray!6, draw=gray!45}
]

\node[block] (guidance) at (0.0,0.0) {LOS / ILOS\\导引律};
\node[core]  (heading)  at (3.3,0.0) {航向整形层\\$\psid \rightarrow \psiref$};
\node[block] (pid)      at (6.7,0.0) {纵向 / 航向\\PID};
\node[plant] (usv)      at (10.0,0.0) {USV\\3DOF 动力学};
\node[core]  (speed)    at (5.0,-2.05) {速度调度层\\$e_s,\tau_{r,\mathrm{need}} \rightarrow \ud$};

\draw[->, thick] (guidance) -- node[above] {$\psid$} (heading);
\draw[->, thick] (heading) -- node[above] {$\psiref$} (pid);
\draw[->, thick] (pid) -- node[above] {$\tau_u,\tau_r$} (usv);

\draw[->, thick] (heading.south) |- node[pos=0.25,left] {$e_s=\psid-\psiref$} (speed.west);
\draw[->, thick, dashed] (pid.south) |- node[pos=0.28,right] {$\tau_{r,\mathrm{need}}$} (speed.east);
\draw[->, thick] (speed.north) -- node[right] {$\ud$} (pid.south);

\draw[->, thick] (usv.south west) |- node[pos=0.30,left] {$(x,y)$} (guidance.south);
\draw[->, thick] (usv.south) |- node[pos=0.32,right] {$r$} (heading.south east);
\draw[->, thick] (usv.south east) |- node[pos=0.28,right] {$u,\psi$} (pid.south east);

\node[note, text width=3.1cm] at (3.3,2.0) {
  \textbf{作用 1}\\
  将几何航向命令变为\\
  满足偏航可达性的\\
  可执行参考
};

\node[note, text width=3.45cm] at (5.0,-3.6) {
  \textbf{作用 2}\\
  根据整形误差和偏航力矩需求\\
  在转角阶段主动降速，\\
  为偏航控制释放推力余量
};

\node[note, text width=3.1cm] at (6.8,2.0) {
  \textbf{保持基线结构}\\
  LOS/ILOS 与 PID 主回路\\
  保持不变，SHCS 作为\\
  可插拔增强层接入
};

\draw[dashed, thick, gray!65, rounded corners=4pt]
  (1.8,1.15) rectangle (6.55,-2.75);
\node[fill=white, inner sep=1pt, font=\small\bfseries, text=orange!70!black]
  at (4.2,1.35) {SHCS 可插拔增强层};

\node[font=\scriptsize, text=gray!70, align=left] at (8.0,-2.05) {虚线表示：\\控制器内部的\\偏航力矩预估支路};

\end{tikzpicture}
\end{document}
